\documentclass[12pt]{article}

\usepackage{amsmath,amssymb}
\usepackage[margin=1in]{geometry}

\begin{document}

\section*{Problem 1 (25\%)}

A box contains 5 red balls, 4 blue balls, and 3 green balls.
Three balls are drawn sequentially without replacement.

Let
\[
A=\{\text{the first ball is red}\},
\]
and let
\[
C=\{\text{at least one of the three drawn balls is red}\}.
\]

Suppose that the only information revealed after the experiment is that
event $C$ occurred; that is, at least one of the three drawn balls was red.

\begin{enumerate}
    \item[(a)] Find $P(C)$.

    \item[(b)] Find $P(A\cap C)$ and hence calculate $P(A\mid C)$.

    \item[(c)] Let $B$ be the event that the second ball is red.
    Calculate $P(B\mid C)$.

    \item[(d)] Calculate $P(A\cap B\mid C)$.

    \item[(e)] Compare $P(A\mid C)$ with $P(A)$.
    Are $A$ and $C$ independent?
    Justify your answer using the definition of independence.
\end{enumerate}

\subsection*{Solution}

\subsubsection*{(a)}

complement of $C$ :none of the three balls is red
\[
P(C^c)
=
\frac{7}{12}\cdot\frac{6}{11}\cdot\frac{5}{10}
=
\frac{7}{44}.
\]

\[
\implies
P(C)
=
1-P(C^c)
=
1-\frac{7}{44}
=
\frac{37}{44}
\]

\subsubsection*{(b)}

\[
\because
A\subseteq C,
\]
\[
\therefore
A\cap C=A.
\]

\[
\implies
P(A\cap C)
=
P(A)
=
\frac{5}{12}
\]

\[
\implies
P(A\mid C)
=
\frac{P(A\cap C)}{P(C)}
=
\frac{\frac{5}{12}}{\frac{37}{44}}
=\frac{55}{111}
\]

\subsubsection*{(c)}

Let
\[
B=\{\text{the second ball is red}\}.
\]

\[
\because
B\subseteq C,
\]
\[
\therefore
P(B\cap C)=P(B).
\]

By symmetry,
\[
P(B)=\frac{5}{12}.
\]

\[
\implies
P(B\mid C)
=
\frac{P(B\cap C)}{P(C)}
=
\frac{\frac{5}{12}}{\frac{37}{44}}.
=
\frac{55}{111}
\]



\subsubsection*{(d)}

\[
\because
P(A\cap B)
=
\frac{5}{12}\cdot\frac{4}{11}
=
\frac{5}{33}.
\]

\[
\therefore
A\cap B\subseteq C,
\]

\[
\therefore
(A\cap B)\cap C=A\cap B.
\]

\[
\implies
P(A\cap B\mid C)
=
\frac{P(A\cap B)}{P(C)}
=
\frac{\frac{5}{33}}{\frac{37}{44}}
=
\frac{20}{111}
\]


\subsubsection*{(e)}


\[
\because
P(A)=\frac{5}{12},
\]
while
\[
P(A\mid C)=\frac{55}{111}.
\]

\[
\therefore
\frac{55}{111}\neq\frac{5}{12},
\]
\[
\implies
P(A\mid C)\neq P(A).
\]

By the definition of independence,
\[
A \text{ and } C \text{ are independent}
\iff
P(A\cap C)=P(A)P(C).
\]

\[
\because
P(A\cap C)=\frac{5}{12},
\]
\[
P(A)P(C)
=
\frac{5}{12}\cdot\frac{37}{44}
=
\frac{185}{528}.
\]


\[
\implies
\frac{5}{12}
=
\frac{220}{528}
\neq
\frac{185}{528},
\]

\[
\therefore
A \text{ and } C \text{ are not independent}
\]

\section*{Problem 2 (25\%)}

\textbf{Multiple Detection System}

A monitoring system uses four independent sensors to detect whether a machine has experienced a failure.
Let $A_i$ denote the event that sensor $i$ raises a false alarm.
The false-alarm probabilities are
\[
P(A_1)=0.10,\qquad
P(A_2)=0.15,\qquad
P(A_3)=0.20,\qquad
P(A_4)=0.25.
\]
Assume that the four events are mutually independent.

The monitoring system issues a warning if at least one sensor raises a false alarm.

\begin{enumerate}
    \item[(a)] Suppose the monitoring system has issued a warning.
    Find the probability that at least two sensors have raised a false alarm, i.e.,
    \[
    P(\text{at least two alarms}\mid A_1\cup A_2\cup A_3\cup A_4).
    \]

    \item[(b)] Suppose now that no independence assumption is made.
    Instead, you are given
    \[
    P(A_i)=0.1,\qquad i=1,2,3,4,
    \]
    and
    \[
    P(A_i\cap A_j)=0.02
    \]
    for every $i\neq j$, together with
    \[
    P(A_i\cap A_j\cap A_k)=0.005
    \]
    for every distinct $i,j,k$, while
    \[
    P(A_1\cap A_2\cap A_3\cap A_4)=0.001.
    \]

    Use the inclusion-exclusion principle to calculate
    \[
    P(A_1\cup A_2\cup A_3\cup A_4).
    \]
\end{enumerate}

\subsection*{Solution}

\begin{enumerate}
    \item[(a)]
    Let
    \[
    W=A_1\cup A_2\cup A_3\cup A_4
    \]
    be the event that the system issues a warning, and let
    \[
    B=\{\text{at least two sensors raise false alarms}\}.
    \]

    Since $B\subseteq W$,
    \[
    P(B\mid W)
    =
    \frac{P(B\cap W)}{P(W)}
    =
    \frac{P(B)}{P(W)}.
    \]

    First,
    \[
    P(W)
    =
    1-P(A_1^c\cap A_2^c\cap A_3^c\cap A_4^c).
    \]

    By independence,
    \[
    \begin{aligned}
    P(W)
    &=
    1-(0.90)(0.85)(0.80)(0.75)\\
    &=
    1-0.459\\
    &=
    0.541.
    \end{aligned}
    \]

    Next,
    \[
    P(B)
    =
    1-P(\text{no alarm})-P(\text{exactly one alarm}).
    \]

    \[
    \therefore
    P(\text{no alarm})
    =
    (0.90)(0.85)(0.80)(0.75)
    =
    0.459.
    \]

    The probability of exactly one alarm is
    \[
    \begin{aligned}
    P(\text{exactly one alarm})
    ={}&
    (0.10)(0.85)(0.80)(0.75)\\
    &+(0.90)(0.15)(0.80)(0.75)\\
    &+(0.90)(0.85)(0.20)(0.75)\\
    &+(0.90)(0.85)(0.80)(0.25)\\
    &=
    0.051+0.081+0.11475+0.153\\
    &=
    0.39975
    \end{aligned}
    \]

    \[
    \therefore
    \begin{aligned}
    P(B)
    &=
    1-0.459-0.39975\\
    &=
    0.14125.
    \end{aligned}
    \]

    \[
    \therefore
    \begin{aligned}
    P(B\mid W)
    &=
    \frac{0.14125}{0.541}\\
    &\approx 0.2611.
    \end{aligned}
    \]

    \[
    \implies
    P(\text{at least two alarms}\mid A_1\cup A_2\cup A_3\cup A_4)
    \approx 0.2611
    \]

    \item[(b)]
    By the inclusion-exclusion principle,
    \[
    \begin{aligned}
    P(A_1\cup A_2\cup A_3\cup A_4)
    ={}&
    \sum_i P(A_i)
    -
    \sum_{i<j}P(A_i\cap A_j)\\
    &+
    \sum_{i<j<k}P(A_i\cap A_j\cap A_k)\\
    &-
    P(A_1\cap A_2\cap A_3\cap A_4).
    \end{aligned}
    \]

    There are
    \[
    \binom{4}{1}=4
    \]
    single events,
    \[
    \binom{4}{2}=6
    \]
    pairwise intersections, and
    \[
    \binom{4}{3}=4
    \]
    triple intersections.

    Therefore,
    \[
    \begin{aligned}
    P(A_1\cup A_2\cup A_3\cup A_4)
    &=
    4(0.1)-6(0.02)+4(0.005)-0.001\\
    &=
    0.4-0.12+0.02-0.001\\
    &=
    0.299.
    \end{aligned}
    \]

    \[
    \implies
    P(A_1\cup A_2\cup A_3\cup A_4)=0.299
    \]

\end{enumerate}


\section*{Problem3}

A factory has three production lines, $L_1$, $L_2$, and $L_3$, which produce
$30\%$, $45\%$, and $25\%$ of the factory's total products, respectively.
The probabilities that a product is defective given that it comes from
$L_1$, $L_2$, or $L_3$ are
\[
P(D\mid L_1)=0.02,\qquad
P(D\mid L_2)=0.03,\qquad
P(D\mid L_3)=0.05,
\]
where $D$ denotes the event that a randomly selected product is defective.

\begin{enumerate}
    \item[(a)] Two products are selected independently from the factory.
    Find the probability that both products are defective.

    \item[(b)] Three products are selected independently from the factory.
    Find the probability that at least one of the three products is defective.

    \item[(c)] Suppose $n$ products are selected independently from the factory.
    Find the smallest integer $n$ such that the probability of obtaining
    at least one defective product is at least $95\%$.
\end{enumerate}

\section*{Solution}

by the law of total probability,
\begin{align*}
P(D)
&=P(D\mid L_1)P(L_1)
 +P(D\mid L_2)P(L_2)
 +P(D\mid L_3)P(L_3)\\
&=(0.02)(0.30)+(0.03)(0.45)+(0.05)(0.25)\\
&=0.006+0.0135+0.0125\\
&=0.032.
\end{align*}

\[
\therefore
P(D)=0.032,
\qquad
P(D^c)=1-0.032=0.968.
\]

\subsection*{(a)}

$\because $ the two products are selected independently
\[
\therefore
P(\text{both defective})
=
P(D)^2
=
(0.032)^2
=
0.001024.
\]

\[
\implies
P(\text{both defective})=0.001024
\]

\subsection*{(b)}
The event ``at least one product is defective'' is the complement of
``none of the three products is defective.''
$\because$ the selections are independent,
\[
P(\text{none defective})
=
(0.968)^3.
\]

\begin{align*}
\therefore
P(\text{at least one defective})
&=1-P(\text{none defective})\\
&=1-(0.968)^3\\
&=0.092960768.
\end{align*}

\[
\implies
P(\text{at least one defective})\approx 0.0930
\]

\subsection*{(c)}

For $n$ independently selected products,
\[
P(\text{no defective products})
=
(0.968)^n.
\]

\[
\therefore
P(\text{at least one defective product})
=
1-(0.968)^n.
\]

\[
\because
1-(0.968)^n\geq 0.95.
\]

\[
\therefore
(0.968)^n\leq 0.05.
\]

Taking natural logarithms,
\[
n\ln(0.968)\leq \ln(0.05).
\]

Since
\[
\ln(0.968)<0,
\]
dividing by $\ln(0.968)$ reverses the inequality:
\[
n\geq
\frac{\ln(0.05)}{\ln(0.968)}.
\]

Hence,
\[
n\geq 92.11\ldots
\]

$\implies$ $n$ must be an integer, the smallest possible value is
\[
n=93
\]



\section*{Problem4}

Two dice are rolled independently. For each die, the probability of showing an even number is \(p\), where
\[
0<p<1.
\]
The two dice have the same distribution, but no assumption is made that the individual outcomes
\(1,2,\ldots,6\) are equally likely.

Define
\[
A=\{\text{the sum of the two dice is even}\},
\]
and
\[
B=\{\text{the sum of the two dice is greater than }8\}.
\]

Determine whether there exists a value of \(p\) such that, for every possible distribution of the individual die outcomes satisfying
\[
P(\text{even})=p,
\]
the events \(A\) and \(B\) are independent.

If such a value of \(p\) exists, find all such values and prove that the independence holds for every admissible distribution.

If no such value exists, prove this by constructing, for each candidate value of \(p\), two admissible distributions that have the same value of \(p\) but give different values of
\[
P(A\cap B)-P(A)P(B).
\]

Your argument should distinguish carefully between independence that is implied by the given information and independence that occurs only for a particular choice of the underlying distribution.

\section*{Solution}

Let the two dice outcomes be \(X\) and \(Y\).
 
\[
P(X\text{ is even})=P(Y\text{ is even})=p,
\]
the two dice are independent
\[
\implies
P(X\text{ is odd})=P(Y\text{ is odd})=1-p.
\]
The event \(A\) occurs if both dice are even or both dice are odd
\[
\implies
P(A)
=
p^2+(1-p)^2.
\]

Thus \(P(A)\) depends only on \(p\).

We now construct two admissible distributions having the same value of \(p\).

\subsection*{Distribution 1}

Let
\[
P(X=1)=1-p,\qquad P(X=2)=p,
\]
and
\[
P(X=k)=0,\qquad k=3,4,5,6.
\]

The same distribution is used for \(Y\).

Clearly,
\[
P(\text{even})=p.
\]

The largest possible sum is
\[
2+2=4,
\]
so
\[
P(B)=0.
\]

\[
\therefore
P(A\cap B)=0,
\]

\[
\therefore
P(A\cap B)-P(A)P(B)
=
0.
\]

Thus, for this particular distribution,
\[
A\text{ and }B\text{ are independent}.
\]

\subsection*{Distribution 2}

Now let
\[
P(X=1)=1-p,\qquad P(X=6)=p,
\]
and
\[
P(X=k)=0,\qquad k=2,3,4,5.
\]

Again, \(Y\) has the same distribution.

\[
\because
P(\text{even})=p.
\]

$\therefore  $  The possible sums are

\[
2,\quad 7,\quad 12.
\]

The event \(B\) occurs only when
\[
X=Y=6.
\]

Therefore
\[
P(B)=p^2.
\]

Also, \(6+6=12\) is even, so
\[
B\subseteq A.
\]

Hence
\[
P(A\cap B)=P(B)=p^2.
\]

Therefore
\begin{align*}
P(A\cap B)-P(A)P(B)
&=
p^2-\left[p^2+(1-p)^2\right]p^2\\
&=
p^2\left[1-p^2-(1-p)^2\right]\\
&=
p^2\left[2p-2p^2\right]\\
&=
2p^3(1-p).
\end{align*}

Since
\[
0<p<1,
\]
we have
\[
2p^3(1-p)>0.
\]

Thus
\[
P(A\cap B)\neq P(A)P(B),
\]
so \(A\) and \(B\) are not independent for this distribution.

The two distributions have the same value of
\[
P(\text{even})=p,
\]
but
\[
P(A\cap B)-P(A)P(B)
=
0
\]
for the first distribution, whereas
\[
P(A\cap B)-P(A)P(B)
=
2p^3(1-p)>0
\]
for the second distribution.

Therefore, for every \(p\in(0,1)\), there exists an admissible distribution for which \(A\) and \(B\) are not independent.

\[
\implies
\text{there is no }p\in(0,1)
\text{ such that }A\text{ and }B
\text{ are independent for every admissible distribution}.
\]

The independence in the first construction occurs only because of that particular underlying distribution; it is not implied by the condition
\[
P(\text{even})=p.
\]


\end{document}